Las experiencias musicales en realidad virtual (VR por sus siglas en inglés) han sido estudiadas principalmente desde formatos cercanos al concierto virtual, donde el interés suele centrarse en la navegación activa, el cambio de perspectiva, el diseño de múltiples vistas, el uso de avatares, la presencia social o las reacciones colectivas de la audiencia \cite{lo2023immersive, lee2025concert, yan2020social, kim2022multiview,  lee2023viewing}. Estos enfoques han permitido ampliar la forma en que una audiencia puede asistir o participar en una presentación musical virtual \cite{onderdijk2023concert}. Sin embargo, dejan abierta una línea de trabajo distinta: el diseño de experiencias musicales en VR donde la canción, y no el concierto, funcione como eje principal de organización visual, narrativa e interactiva.

A partir de esta línea de trabajo, el proyecto se justifica en la necesidad de estudiar cómo una canción puede estructurar una experiencia inmersiva sin depender de las lógicas propias de un concierto virtual. Mientras que los conciertos VR suelen organizarse alrededor del escenario, la audiencia, el artista o la interacción social, un videoclip musical VR puede tomar como punto de partida la progresión temporal de la obra. Esto permite que la experiencia se diseñe desde los momentos expresivos de la canción y no desde la reproducción de una presentación en vivo.

Por esta razón, se plantea una secuenciación manual de eventos, en la que las activaciones visuales, sonoras e interactivas se diseñan de acuerdo con la identidad artística y la estructura de una canción específica. Esta decisión se justifica porque permite conservar una relación directa entre cada estímulo y la intención narrativa o expresiva del recorrido musical. De esta forma, los cambios visuales, elementos sonoros, transiciones de ambiente e interacciones no se integran únicamente como recursos técnicos, sino como componentes que acompañan momentos específicos de la canción y refuerzan la coherencia de la experiencia.

La propuesta también se diferencia de experiencias interactivas basadas principalmente en generación procedural. Aunque estos métodos permiten producir contenido dinámico mediante reglas algorítmicas, distintos autores señalan la importancia de conservar criterios de diseño, objetivos expresivos y control autoral sobre la experiencia resultante \cite{yannakakis2011experience,khaled2013design, craveirinha2015designing}. Esta diferencia resulta relevante porque el propósito del proyecto no es automatizar la creación de contenido, sino diseñar una experiencia que responda a una obra musical concreta. Por ello, la secuenciación manual permite mantener el control sobre la relación entre música, imagen, interacción y narrativa.

Desde el punto de vista técnico, la Arquitectura Dirigida por Eventos (EDA por sus siglas en inglés) ofrece una base pertinente para organizar la experiencia porque permite estructurar el sistema mediante eventos, productores, consumidores y mecanismos de gestión, separando los componentes que generan señales de los elementos que reaccionan a ellas \cite{oberhauser2023vr}. En este proyecto, EDA se justifica porque la experiencia requiere coordinar múltiples elementos del entorno virtual sin perder su relación con la estructura temporal de la canción. Además, investigaciones previas han utilizado enfoques basados en eventos en entornos virtuales para representar flujos de información y coordinar interacciones multisensoriales mediante eventos discretos \cite{oberhauser2023vr, edmunds2006mixed}. Por lo tanto, esta arquitectura permite que la secuenciación diseñada pueda ejecutarse de forma organizada y coherente dentro del entorno VR.

La evaluación mediante el I-Group Presence Questionnaire (IPQ) permite analizar la experiencia desde la percepción del usuario y no únicamente desde el funcionamiento técnico del sistema. Este instrumento considera dimensiones como presencia espacial, involucramiento y realismo, las cuales permiten estudiar cómo el usuario percibe su relación con el entorno virtual \cite{berkman2021igroup}. Además, estudios sobre conciertos VR han relacionado la inmersión, la presencia y las formas de interacción con la satisfacción, el flujo y la intención de participar en futuras experiencias virtuales \cite{onderdijk2023concert,lo2023immersive, lee2023viewing}. Por esta razón, el IPQ permite valorar si la secuenciación inmersiva basada en eventos favorece la sensación de presencia dentro del videoclip musical en VR.

En el contexto de artistas guatemaltecos, el proyecto plantea un formato alternativo para presentar una canción mediante una experiencia visual, auditiva e interactiva. Las audiencias de conciertos VR han valorado la singularidad de la experiencia y la posibilidad de interactuar con visuales y entornos que no serían posibles en un escenario físico tradicional \cite{onderdijk2023concert}. De esta manera, el proyecto adquiere relevancia al proponer una experiencia que integra música, diseño inmersivo y arquitectura de software para explorar nuevas formas de presentación artística.
